% ============================================================
% BTH04 - Yeu cau 4: Phan tich 3 du an nguon mo voi 3 giay phep khac nhau
% Overleaf: New Project -> Upload Project (zip). Compiler: pdfLaTeX.
% ============================================================
\documentclass[runningheads]{llncs}

\usepackage{iftex}
\ifPDFTeX
  \usepackage[T5]{fontenc}
  \usepackage[utf8]{inputenc}
\else
  \usepackage{fontspec}
  \setmainfont{TeX Gyre Termes}
\fi
\usepackage[vietnamese]{babel}
\usepackage{tabularx}
\usepackage{booktabs}
\usepackage{array}
\usepackage[unicode]{hyperref}
\usepackage{color}
\renewcommand\UrlFont{\color{blue}\rmfamily}
\newcolumntype{L}{>{\raggedright\arraybackslash}X}

\begin{document}

\title{So sánh giấy phép GPL, MIT và Apache 2.0 qua ba dự án nguồn mở: GNU Bash, React và TensorFlow}
\titlerunning{So sánh giấy phép nguồn mở}
\author{Phạm Lê Trường Thịnh\inst{1}\orcidID{0009-0002-5606-6810}}
\authorrunning{Phạm Lê Trường Thịnh}
\institute{Trường Công nghệ Thông tin và Truyền thông, Đại học Cần Thơ, Việt Nam\\
\email{email@student.ctu.edu.vn}}
\maketitle

\begin{abstract}
Giấy phép quyết định ai được dùng, sửa và phân phối lại một phần mềm nguồn mở, và với điều
kiện gì. Báo cáo phân tích ba dự án tiêu biểu cho ba họ giấy phép: GNU Bash (GPLv3, copyleft
mạnh), React (MIT, dễ dãi) và TensorFlow (Apache 2.0, dễ dãi có điều khoản bằng sáng chế).
Chúng tôi so sánh quyền sử dụng, chỉnh sửa, phân phối, tác động tới cộng đồng và khả năng
thương mại hóa, từ đó đề xuất loại giấy phép phù hợp cho từng loại dự án, kể cả dự án
\texttt{student-performance-ml} của chính tác giả.
\keywords{Giấy phép nguồn mở \and GPL \and MIT \and Apache 2.0 \and Bản quyền nội dung số.}
\end{abstract}

\section{Giới thiệu}\label{sec:intro}
Theo định nghĩa của Open Source Initiative, phần mềm nguồn mở phải cho phép tự do sử dụng,
nghiên cứu, sửa đổi và phân phối lại~\cite{osd}. Các giấy phép chia thành hai nhóm chính:
\emph{copyleft} (ví dụ GPL) bắt buộc bản phái sinh giữ nguyên giấy phép, và \emph{dễ dãi}
(permissive, ví dụ MIT, Apache 2.0) chỉ yêu cầu giữ thông báo bản quyền. Phần~\ref{sec:projects}
phân tích ba dự án, Phần~\ref{sec:compare} so sánh và Phần~\ref{sec:recommend} đưa ra khuyến nghị.

\section{Phân tích ba dự án}\label{sec:projects}

\subsection{GNU Bash – GPLv3}\label{sec:bash}
\textbf{Mục tiêu và phạm vi.} Bash (Bourne-Again SHell) là trình thông dịch dòng lệnh của
dự án GNU, là shell mặc định trên phần lớn bản phân phối Linux và được dùng rộng rãi để viết
script tự động hóa hệ thống~\cite{bash}.\\
\textbf{Cộng đồng.} Bash được bảo trì chủ yếu bởi một người bảo trì chính (Chet Ramey) với
đóng góp qua danh sách thư của GNU. Cộng đồng người dùng rất lớn, nhưng nhóm phát triển lõi nhỏ.\\
\textbf{Giấy phép.} Từ phiên bản 4.0, Bash phát hành theo GPL phiên bản 3 hoặc mới hơn~\cite{gpl3}.
Người dùng được chạy, sửa và phân phối tự do; nhưng nếu phân phối bản đã sửa (hoặc phần mềm
liên kết với nó) thì phải công bố mã nguồn theo cùng giấy phép GPL. GPLv3 còn cấm \emph{tivoization}
(khóa phần cứng để ngăn chạy bản đã sửa) và có điều khoản cấp phép bằng sáng chế.

\subsection{React – MIT}\label{sec:react}
\textbf{Mục tiêu và phạm vi.} React là thư viện JavaScript để xây dựng giao diện người dùng
theo mô hình thành phần (component), do Meta phát triển, dùng cho web và (qua React Native) cho
ứng dụng di động~\cite{react}.\\
\textbf{Cộng đồng.} React là một trong những kho mã phổ biến nhất trên GitHub, với hàng nghìn
người đóng góp và hệ sinh thái thư viện khổng lồ trên npm.\\
\textbf{Giấy phép.} React dùng giấy phép MIT từ năm 2017 (trước đó là BSD kèm điều khoản bằng
sáng chế gây tranh cãi)~\cite{mit}. MIT chỉ yêu cầu giữ thông báo bản quyền và giấy phép trong
mọi bản sao; được dùng trong phần mềm đóng, thương mại mà không phải công bố mã nguồn; không có
bảo hành và không đề cập bằng sáng chế.

\subsection{TensorFlow – Apache 2.0}\label{sec:tf}
\textbf{Mục tiêu và phạm vi.} TensorFlow là nền tảng học máy mã nguồn mở do Google phát triển,
hỗ trợ huấn luyện và triển khai mô hình trên CPU, GPU, TPU, thiết bị di động và trình duyệt~\cite{tf}.\\
\textbf{Cộng đồng.} Dự án có hàng nghìn người đóng góp, được dùng rộng rãi trong nghiên cứu và
công nghiệp, quy trình đóng góp qua GitHub có kiểm duyệt của Google và yêu cầu ký thỏa thuận CLA.\\
\textbf{Giấy phép.} Apache 2.0~\cite{apache} là giấy phép dễ dãi như MIT nhưng chi tiết hơn: cấp
rõ quyền sử dụng bằng sáng chế từ người đóng góp, tự động chấm dứt quyền này nếu người dùng kiện
về bằng sáng chế, yêu cầu ghi chú các tệp đã sửa đổi và giữ tệp \texttt{NOTICE}; không cấp quyền
dùng nhãn hiệu.

\section{So sánh}\label{sec:compare}
Bảng~\ref{tab:compare} tổng hợp khác biệt giữa ba giấy phép.

\begin{table}[t]
\centering
\caption{So sánh GPLv3, MIT và Apache 2.0.}\label{tab:compare}
\small
\begin{tabularx}{\linewidth}{@{}p{2.3cm}LLL@{}}
\toprule
Tiêu chí & GPLv3 (Bash) & MIT (React) & Apache 2.0 (TensorFlow) \\
\midrule
Sử dụng & Tự do, mọi mục đích & Tự do, mọi mục đích & Tự do, mọi mục đích \\
Chỉnh sửa & Được & Được & Được, phải ghi chú tệp đã sửa \\
Phân phối lại & Phải kèm mã nguồn, giữ GPL & Chỉ cần giữ thông báo bản quyền & Giữ giấy phép và tệp NOTICE \\
Bằng sáng chế & Có cấp quyền & Không đề cập & Cấp rõ, có điều khoản trả đũa \\
Tích hợp vào sản phẩm đóng & Không (bản phái sinh phải mở) & Có & Có \\
Tương thích & Nhận mã MIT/Apache; không cho ngược lại & Gần như mọi giấy phép & Tương thích GPLv3, không tương thích GPLv2 \\
Tác động cộng đồng & Giữ mọi cải tiến ở dạng mở & Dễ được doanh nghiệp chấp nhận & Doanh nghiệp yên tâm về pháp lý \\
\bottomrule
\end{tabularx}
\end{table}

\subsection{Quyền sử dụng, chỉnh sửa, phân phối}
Cả ba giấy phép đều cho phép dùng và sửa tự do. Khác biệt nằm ở \emph{phân phối}: GPL lan truyền
nghĩa vụ mở mã sang bản phái sinh, còn MIT và Apache thì không (Bảng~\ref{tab:compare}).

\subsection{Tác động đến cộng đồng phát triển}
GPL bảo đảm mọi cải tiến quay lại cộng đồng, phù hợp với triết lý phần mềm tự do của GNU. MIT
có rào cản thấp nhất nên React được tích hợp vào vô số sản phẩm thương mại, kéo theo nhiều người
đóng góp. Apache 2.0 giúp các công ty lớn đóng góp mà không lo rủi ro bằng sáng chế, đó là lý do
nhiều dự án hạ tầng và AI chọn nó.

\subsection{Thương mại hóa và tích hợp vào sản phẩm đóng}
Doanh nghiệp có thể nhúng React hoặc TensorFlow vào phần mềm độc quyền. Với Bash, doanh nghiệp vẫn
có thể \emph{dùng} hoặc gọi Bash như một chương trình riêng, nhưng nếu sửa và phân phối Bash thì
phải công bố mã nguồn phần đã sửa.

\section{Nhận xét và khuyến nghị}\label{sec:recommend}
\begin{itemize}
  \item \textbf{GPLv3}: phù hợp với công cụ hệ thống, phần mềm cộng đồng muốn ngăn bị ``đóng''
        lại (như Bash ở Mục~\ref{sec:bash}).
  \item \textbf{MIT}: phù hợp với thư viện nhỏ, mã nguồn học tập, dự án muốn lan tỏa tối đa
        (như React ở Mục~\ref{sec:react}).
  \item \textbf{Apache 2.0}: phù hợp với dự án lớn có nhiều doanh nghiệp đóng góp hoặc có rủi ro
        bằng sáng chế, ví dụ nền tảng AI (TensorFlow ở Mục~\ref{sec:tf}).
\end{itemize}
Với dự án \texttt{student-performance-ml} của tác giả (mã phân tích dữ liệu cho mục đích học tập),
giấy phép \textbf{MIT} được chọn vì đơn giản và cho phép người khác tái sử dụng tự do. Dữ liệu đi
kèm vẫn giữ giấy phép CC BY 4.0 của tác giả gốc, nên README ghi rõ hai giấy phép tách biệt.

\begin{thebibliography}{8}
\bibitem{apache} Apache Software Foundation: Apache License, Version 2.0 (2004),
  \url{https://www.apache.org/licenses/LICENSE-2.0}
\bibitem{bash} Free Software Foundation: GNU Bash, \url{https://www.gnu.org/software/bash/}
\bibitem{gpl3} Free Software Foundation: GNU General Public License, Version 3 (2007),
  \url{https://www.gnu.org/licenses/gpl-3.0.html}
\bibitem{mit} Open Source Initiative: The MIT License, \url{https://opensource.org/license/mit}
\bibitem{osd} Open Source Initiative: The Open Source Definition,
  \url{https://opensource.org/osd}
\bibitem{react} Meta Open Source: React – The library for web and native user interfaces,
  \url{https://github.com/facebook/react}
\bibitem{tf} Abadi, M., et al.: TensorFlow: Large-scale machine learning on heterogeneous
  distributed systems. arXiv:1603.04467 (2016). \url{https://doi.org/10.48550/arXiv.1603.04467}
\end{thebibliography}

\end{document}
